\chapter{Problem Definition}

\section{Problem Statement}
Organizations keep their legal records (contracts, NDAs, leases, employment letters, affidavits)
as a mix of digital files, scanned pages and handwritten stamp-paper documents, stored in folders that
are not connected to each other. As a result:
\begin{enumerate}
\item \textbf{Facts are slow to find.} Answering \emph{``When was the NDA with Acme signed?''} or
\emph{``Which leases expire this year?''} requires opening and reading files manually, and questions
across many documents are practically impossible to answer.
\item \textbf{Scanned and handwritten documents are invisible to search.} Their content exists only as
images.
\item \textbf{Reference material is hard to use.} Finding the relevant provision in a 100-page law book
such as the Contract Act 1872 takes time and expertise.
\item \textbf{Access is uncontrolled.} Shared folders are either open to everyone or require manual
sharing, and there is no record of who opened or changed a document.
\item \textbf{Deadlines are missed.} Expiry dates and notice periods are buried inside documents, so
agreements renew or lapse without anyone noticing.
\end{enumerate}
General AI document assistants reduce the reading effort but do not solve the problem for an
organization: they cannot prove where an answer came from at the level of a single field, they do not
keep a structured register of records, and they do not enforce who may see what.

\section{Proposed Solution Overview}
\sysname{} is a shared, access-controlled legal records system for one organization at a time. Every
uploaded document passes through a processing pipeline that reads all pages (including scanned and
handwritten ones), classifies the document as a \emph{record} or \emph{reference} material, splits it
along its section structure, indexes it for keyword and semantic search, and, for records, extracts
structured fields with page-level evidence. Low-confidence or unverified values go to a review queue
where a manager confirms or corrects them.

Employees ask questions in plain English. A query router decides whether the question asks for a
stored fact (answered from the records table), an explanation from document text (answered from
retrieved passages with citations), or a set of documents (answered with a ranked list). Every answer
links to the source page, and the document viewer highlights the evidence. Access rules are applied
before any search, and query history, document versions, field edits and all user actions are kept.

\section{Objectives of the Proposed System}
\begin{enumerate}
\item Provide one shared document repository for an organization with role-based, department-based and
named-person access control.
\item Read digital, scanned and handwritten (English) documents and make all of them searchable.
\item Automatically distinguish records from reference material and extract eleven structured fields
from each record, each backed by a verbatim quote and a page number.
\item Answer fact questions (e.g.\ signing dates) directly from the extracted records, and explanation
questions (e.g.\ from a 100-page law book) from retrieved passages with citations.
\item Open the source document at the cited page with the evidence highlighted.
\item Maintain four types of history: query history, document versions, field-edit history and an
audit log.
\item Remind responsible people before agreements expire or notice deadlines pass.
\item Fine-tune a clause-classification model and evaluate every AI component quantitatively on public
legal datasets.
\end{enumerate}

\section{Scope of the System}
\textbf{In scope:}
\begin{itemize}
\item Organization registration, member and department management, four roles (Admin, Manager,
Employee, Auditor).
\item Upload of PDF, DOCX, PNG and JPG documents; new versions of an existing document and comparison
between versions.
\item Text extraction from digital PDFs and DOCX files, OCR of printed scans, recognition of English
handwriting.
\item Record/reference classification, structure-aware chunking, embeddings, clause labelling and record
field extraction with grounding.
\item Records register, review queue, verification and correction of fields.
\item Assistant with lookup, knowledge and search answers, citations and conversation history.
\item Audit log, in-app notifications and expiry / notice reminders.
\item Fine-tuning and evaluation of the clause classifier; evaluation of retrieval, extraction, routing,
answers and handwriting.
\end{itemize}
\textbf{Out of scope:} Urdu or other non-English documents, drafting or e-signing of contracts, legal
advice, native mobile applications, and multi-organization membership for one user.

\section{System Architecture and Modules}
\sysname{} follows a three-tier client-server architecture with a separate background worker
(Figure~\ref{fig:highlevel}). The Vue.js single-page application communicates with the Django REST
API. Uploads are stored and queued; the worker processes them asynchronously, so the user interface
never waits for OCR or model calls. PostgreSQL with pgvector stores relational data, embeddings and
full-text indexes together. External LLM services are reached only through one gateway module.

\begin{figure}[H]\centering
\begin{tikzpicture}[node distance=6mm and 6mm]
\node[box,minimum width=11cm,fill=gray!8] (ui) {\textbf{Vue.js Single-Page Application}\\ Assistant $\cdot$ Documents $\cdot$ Records $\cdot$ Review $\cdot$ Team $\cdot$ Audit};
\node[box,below=10mm of ui,minimum width=11cm] (api) {\textbf{Django REST API}\\ JWT authentication $\cdot$ RBAC $\cdot$ throttling $\cdot$ audit};
\node[box,below=10mm of api,text width=34mm] (m4) {M4 Legal\\ Assistant};
\node[box,left=3mm of m4,text width=34mm] (m1) {M1 Organization \&\\ Access};
\node[box,right=3mm of m4,text width=34mm] (m3) {M3 Records \&\\ Verification};
\node[box,below=4mm of m1,text width=34mm] (m5) {M5 History, Audit\\ \& Notifications};
\node[box,below=4mm of m4,text width=34mm] (m2) {M2 Document Ingestion\\ (background worker)};
\node[box,below=4mm of m3,text width=34mm] (m6) {M6 Model Training\\ \& Evaluation};
\node[db,below=12mm of m2] (db) {PostgreSQL 16\\ + pgvector};
\node[box,fill=gray!8,text width=34mm] (llm) at (m6 |- db) {External LLM / VLM\\ (Groq, Gemini)};
\draw[<->,thick] (ui) -- node[lbl,right]{HTTPS / JSON} (api);
\draw[->,thick] (api) -- (m4.north);
\draw[->,thick] (api.south) -| (m1.north);
\draw[->,thick] (api.south) -| (m3.north);
\draw[<->,thick] (m2) -- (db);
\draw[<->,thick] (m5.south) |- (db.west);
\draw[->,dashed,thick] (m2.south east) -- (llm.north west);
\draw[->,dashed,thick] (m6) -- (llm);
\end{tikzpicture}
\caption{High-Level System Architecture of \sysname{}}\label{fig:highlevel}
\end{figure}

\subsection{Module 1: Organization and Access Management}
Registers an organization and its first administrator, authenticates users with JSON Web Tokens, and
lets administrators add members, assign roles and departments. It defines the access model used by
every other module: a user sees all organization-wide documents, the documents of their department,
documents they uploaded and documents shared with them by name; Admins and Auditors see all documents
of the organization.

\subsection{Module 2: Document Ingestion and Processing}
Validates uploads, stores each file as a numbered version, and queues it for the background worker.
The worker extracts text page by page, chooses between the digital text layer, Tesseract OCR and the
vision-language model for each page, classifies the document, chunks it by section, computes
embeddings, labels contract clauses with the fine-tuned classifier, and extracts record fields.
Users can upload new versions and compare any two versions line by line.

\subsection{Module 3: Records Register and Verification}
Presents all records of the organization as one register (title, type, parties, dates, status) that
can be filtered and sorted. Each field shows its evidence quote and page. Fields marked
\emph{Needs review} appear in a review queue where Managers and Admins verify or correct them; every
change is stored in the field's edit history.

\subsection{Module 4: Legal Assistant}
Accepts a question, optionally limited to selected documents. A router (an LLM with a keyword fallback)
classifies the question as LOOKUP, KNOWLEDGE or SEARCH and extracts filters such as document type,
party, field and date range. Lookups are answered from the records table; knowledge questions are
answered by an LLM from hybrid-retrieved passages with numbered citations; searches return ranked
documents. Each conversation is saved.

\subsection{Module 5: History, Audit and Notifications}
Writes an append-only audit event for logins, uploads, views, downloads, searches, field edits, access
changes and membership changes. It also produces notifications: when a document is processed or fails,
when fields need review, and reminders 90, 30 and 7 days before (and on the day of) an expiry or a
notice deadline, sent only to people who can see the document.

\subsection{Module 6: Model Training and Evaluation}
Contains the scripts that fine-tune the clause classifier on LEDGAR, the TF-IDF and zero-shot LLM
baselines, and the evaluation suite that measures each AI component on CUAD, LEDGAR, a labelled
routing set and a handwriting set.

\section{Assumptions and Dependencies}
\begin{itemize}
\item Users have a modern web browser and network access to the server.
\item Uploaded documents are in English and reasonably legible.
\item An API key for at least one supported LLM provider (Groq, Gemini, OpenAI or a local Ollama server)
is configured; without it, upload, search, lookup and access control still work, while generated
answers and field extraction are disabled.
\item The server provides PostgreSQL 16 with pgvector and the Tesseract OCR engine.
\item The organization assigns roles responsibly; an administrator is trusted to manage members.
\end{itemize}

\section{Chapter Summary}
This chapter defined the problem of finding, trusting and protecting information in an organization's
legal documents, presented the proposed solution, its objectives and scope, and described the six
modules of the system. The next chapter analyzes the requirements in detail.
